\chapter{Running a Currency-Free Economy}
\label{ch:gm-currency-free}

\begin{tcolorbox}[colback=blue!5!white,colframe=blue!70!black,title=GM's Note: The Ledger Is Fiction]
\emph{“You don’t need to track every copper. Track what matters: the safehouse the players bled for, the oath they swore to the pirate queen, the favor they owe the fox spirit. That’s the real economy.”}
\end{tcolorbox}

This chapter provides guidance for running the currency-free economy described in the Player’s Guide (\ref{ch:currency-free-economy}). As the GM, your role is to translate player investments—XP, barter, Strings, Obligation, and Trusts—into meaningful story outcomes. You are not an accountant; you are a curator of consequence.

\section{Core Principles for the GM}
\label{sec:gm-currency-principles}

\begin{enumerate}
  \item \textbf{Hand-wave the trivial.} If a player asks for something small and sensible (a meal, a horse, a basic tool), say yes. Scarcity is a tool, not a default.
  \item \textbf{Significant gains cost significant resources.} Use the XP cost table (\ref{subsec:xp-cost-items}) as a baseline, but be flexible: barter, Strings, or Obligation are equally valid.
  \item \textbf{Debt is a story engine.} When players accept a Debt Clock, they are inviting future complications. Don’t let debts linger – call them in at dramatically appropriate moments.
  \item \textbf{Trusts are player-driven factions.} Encourage players to form Trusts, then use them to generate missions, rivals, and hard choices.
  \item \textbf{XP spent is XP remembered.} When a player spends XP on an asset or a follower, make that asset matter in the fiction. A safehouse should be threatened; a follower should have opinions.
\end{enumerate}

\section{Adjudicating Barter and Strings}
\label{sec:gm-barter-strings}

\subsection{When to Allow Barter}
\label{subsec:when-barter}

Barter is appropriate when the player character has something the other party genuinely wants. Ask yourself:

\begin{itemize}
  \item Does the PC possess a unique item, secret, or service?
  \item Does the NPC have a clear motivation (greed, fear, ambition)?
  \item Is the exchange roughly equal in narrative weight?
\end{itemize}

If the answer to all three is yes, allow barter. If not, the NPC may demand XP, Obligation, or a String instead.

\paragraph{Example:} Sera offers a stolen Ashaani spirit-lamp to a Sidhi smuggler in exchange for passage on her ship. The smuggler wants the lamp (it can hide her cargo from inspectors). The exchange is balanced. No roll is needed unless the GM wants to introduce tension (e.g., a Deception roll to hide the lamp’s curse).

\subsection{Rolling for Barter}
\label{subsec:rolling-barter}

For contentious or high-stakes barter, call for a roll:

\begin{itemize}
  \item \textbf{Skill:} Presence+Sway (negotiation), Wits+Streetwise (knowing the right price), or Spirit+Lore (appealing to shared customs).
  \item \textbf{DV:} 2 (friendly counterparty) to 5 (desperate or hostile counterparty).
  \item \textbf{Position:} Controlled for fair trade; Desperate if the PC is under pressure.
\end{itemize}

On a \textbf{Clean Success}, the barter succeeds as proposed. On a \textbf{Success with SB}, the trade happens but the GM gains 1 SB to spend on a complication (e.g., the item is stolen goods, the NPC will betray them later). On a \textbf{Partial}, the trade happens but the PC must add a sweetener (a small favor, a promise, 1 XP). On a \textbf{Miss}, the NPC refuses; the GM may offer a Devil’s Bargain (accept a Debt Clock to salvage the deal).

\subsection{Spending Strings}
\label{subsec:spending-strings-gm}

When a player wants to spend a String (see \ref{subsec:strings-currency}), ask:

\begin{itemize}
  \item \textbf{Is the String relevant?} A “Minor String” from a grain inspector won’t help negotiate with a pirate queen. If the String is a poor fit, the player can still spend it, but reduce its effective value by one tier (e.g., Minor → trivial, only a +1 die bonus).
  \item \textbf{Does the String need a roll?} Usually not – a String represents an established relationship. However, if the NPC is hostile or the request is extreme, call for a roll (Presence+Sway, DV 3). On a failure, the String is still consumed, but the NPC demands something extra.
  \item \textbf{What happens when a String is spent?} The narrative permission is used up. The NPC may say “We’re even now.” The String is removed from the player’s sheet. If the player wants a similar favor later, they must earn a new String.
\end{itemize}

\subsection{Introducing New Strings}
\label{subsec:introducing-strings}

Strings are a great reward for roleplaying. Award a Minor String when a player:
\begin{itemize}
  \item Saves an NPC’s life or reputation.
  \item Completes a significant side quest for a faction.
  \item Reveals a secret that benefits the NPC.
  \item Pays a Debt Clock segment through service (see \ref{sec:gm-debt-clocks}).
\end{itemize}

Award a Standard or Major String for campaign-level deeds: brokering a peace treaty, exposing a conspiracy, or retrieving a legendary artifact.

\section{Managing Debt Clocks}
\label{sec:gm-debt-clocks}

Debt Clocks are the currency-free equivalent of loans with interest. They are story engines, not punishments.

\subsection{Creating a Debt Clock}
\label{subsec:creating-debt-clock}

When a player accepts a Debt Clock (e.g., in exchange for an asset or a favor), create a visible clock (4, 6, or 8 segments) with a clear label: \textit{“Debt to the Vermilion Cord [6]”} or \textit{“Oath to the Fox Spirit [4]”}.

\subsection{Ticking the Debt Clock}
\label{subsec:ticking-debt-clock}

Advance the clock by 1 segment when:
\begin{itemize}
  \item The player fails a roll related to the debt (e.g., failing to deliver cargo on time).
  \item The GM spends 2 Story Beats (SB) on a complication tied to the debt.
  \item A session ends without the player making progress toward repayment (tick once per session).
\end{itemize}

When the clock reaches its final segment, the debt is \textbf{due}. The creditor makes a demand. The demand should be significant but not impossible:
\begin{itemize}
  \item A dangerous mission (steal from a rival, assassinate a target).
  \item A public declaration of loyalty (which may alienate another faction).
  \item A sacrifice of equal value (a personal Asset, a magical item, a String).
\end{itemize}
If the player refuses, the creditor becomes an enemy (advance the Crisis clock, start a new Rival clock, or send enforcers).

\subsection{Clearing Debt Clock Segments}
\label{subsec:clearing-debt-segments}

Players can clear segments by:
\begin{itemize}
  \item Spending 2 XP per segment (treat as buying off the debt with hard work).
  \item Performing a service (GM discretion – usually a scene or a roll clears 1–2 segments).
  \item Offering an equivalent String or Asset to the creditor.
\end{itemize}

\subsection{Debt as a Reward}
\label{subsec:debt-reward}

Do not treat Debt Clocks as purely negative. They give players access to resources they couldn’t otherwise afford. A party that accepts a Debt Clock [6] for a ship gains a powerful asset immediately – the clock is the price of adventure. Make the eventual demand exciting, not punitive.

\section{Running Trusts (Player-Owned Patrons)}
\label{sec:gm-trusts}

Trusts are player-led factions. They give players a sense of ownership and generate endless plot hooks.

\subsection{Helping Players Form a Trust}
\label{subsec:helping-trust}

When players express interest in forming a Trust (see \ref{sec:forming-trust}), guide them through:
\begin{itemize}
  \item \textbf{Naming the Trust} – should reflect purpose and tone.
  \item \textbf{Defining its Seat} – a physical headquarters (warehouse, guild hall, ship).
  \item \textbf{Choosing a Purpose} – smuggling, trade, information, mercenary work, etc.
  \item \textbf{Establishing an initial Debt Clock} – the cost of founding. Clear it through a short introductory mission.
\end{itemize}
Require the party to contribute 10 XP from their pool (see \ref{subsec:creating-pool}). If they don’t have enough, let them work toward it over a few sessions.

\subsection{Trust Obligation}
\label{subsec:trust-obligation-gm}

Members incur Obligation to the Trust when they ask for its help (borrowing an Asset, requesting a Follower, etc.). Track this separately from other Obligation. When a member’s Obligation exceeds their Spirit+Presence, the Trust demands a task. This is your chance to introduce a mission that advances the Trust’s goals – and puts the party in conflict with other factions.

\paragraph{Sample Trust Demands:}
\begin{itemize}
  \item “The Silkhand League needs you to eliminate a rival merchant’s shipment.”
  \item “Our spy network has gone silent in Port Shed. Investigate.”
  \item “The harbor master is squeezing our dues. Find leverage on him.”
  \item “One of our own has been kidnapped by the Veiled Sisters. Rescue them.”
\end{itemize}
If a player refuses a Trust demand, they are expelled and lose all Trust benefits. Expulsion should be a major story beat – the character becomes a pariah among their own creation.

\subsection{Growing the Trust}
\label{subsec:growing-trust-gm}

As players invest more XP into the Trust (see \ref{subsec:trust-growth}), expand its capabilities and its enemies. A Trust with a spy network (Standard Asset) attracts the attention of rival intelligence agencies. A Trust with a fortress (Major Asset) becomes a target for sieges or political intrigue. Use the Trust’s growth to escalate the campaign.

\subsection{Trusts as Antagonists}
\label{subsec:trust-antagonists}

A Trust can also be an enemy – a rival guild, a corrupt merchant house, or a fanatical cult. Treat it as a faction with its own clocks (see \ref{sec:clocks}) and Strings. Players might infiltrate it, steal from it, or try to destroy it. The rules for Trusts are symmetrical: you can run an NPC Trust using the same mechanics.

\section{Handling Common Player Behaviors}
\label{sec:common-player-behaviors}

\subsection{Hoarding XP}
\label{subsec:hoarding-xp}

Players may be reluctant to spend XP on non-personal advancement. Remind them:
\begin{itemize}
  \item XP spent on assets and Trusts is not lost – it creates durable story elements that benefit the whole party.
  \item Unspent XP does nothing. A safehouse or a follower is tangible power.
  \item The game’s pacing assumes some XP will go toward external resources. If players hoard, they may lack leverage in faction-level conflicts.
\end{itemize}
If hoarding persists, introduce situations where personal power alone is insufficient: a locked vault that requires a spy network (Asset), a siege that needs a safehouse, a negotiation that demands a Trust’s backing.

\subsection{Farming Boons for XP Conversion}
\label{subsec:farming-boons}

In the currency-free economy, players can still convert 2 Boons → 1 XP (see \ref{sec:boons}). If they try to farm Boons by intentionally failing trivial rolls, invoke the anti-fishing measures (\ref{subsec:boons-anti-fishing}): no Boons from safe, repetitive, or stakes-free actions. Remind them that Boons are for dramatic moments, not XP grinding.

\subsection{Abusing Strings}
\label{subsec:abusing-strings}

If players try to use the same String repeatedly (“I call in the grain inspector’s favor again”), rule that a String is exhausted after one meaningful use. The NPC may say “We’re even now.” Players can earn a new String by doing another favor.

\subsection{Overloading the Party Pool}
\label{subsec:overloading-pool}

If players contribute all their XP to the party pool and neglect personal advancement, they may become underpowered in direct challenges. Warn them gently: “Your character hasn’t improved their combat skills in a while – are you sure you want to put everything into the Trust?” Offer a session or two of personal advancement before the next big mission.

\section{Integrating with Existing Systems}
\label{sec:integrating-systems}

\subsection{Supply Clock}
\label{subsec:supply-currency}

The Supply Clock (see \ref{sec:supply-clock}) represents scarcity of mundane resources. Use it when the party is isolated, pursued, or in a hostile environment. Do not use it to nickel-and-dime players over everyday costs – that defeats the purpose of the currency-free system.

\subsection{Assets and Followers}
\label{subsec:assets-followers-currency}

The XP costs for Assets and Followers remain unchanged. The only difference is that players never pay coin for upkeep – only XP or downtime actions (see \ref{subsec:upkeep-options}). If a player wants to “buy” an Asset with barter or Strings, use the equivalencies in \ref{subsec:strings-currency}.

\subsection{Terrestrial Patrons}
\label{subsec:terrestrial-currency}

A terrestrial patron (see \ref{sec:terrestrial-patrons}) is a natural fit for the currency-free economy. Instead of paying coin for a patron’s favor, the party accepts a Debt Clock. The patron’s demands become adventure hooks. Treat Trusts as a special case of terrestrial patron – one controlled by the players themselves.

\subsection{Magic and Components}
\label{subsec:magic-components}

Magical components (rare herbs, spirit dust, etc.) should be treated as barter goods or Strings, not coin. A player might acquire a “dragon’s tooth” (String value: Major) by completing a quest, then use it to power a High Rite without spending XP. Alternatively, they can spend XP directly to represent the effort of finding or crafting the component.

\section{GM Tools and Quick Reference}
\label{sec:gm-currency-tools}

\subsection{At-a-Glance: When to Charge What}
\label{subsec:charge-what}

| Player wants... | Charge... |
|----------------|-----------|
| A meal, a horse, a basic tool | Nothing (hand-wave) |
| A safehouse, a spy network, a workshop | XP (4–12) or barter/String of equal value |
| A follower (Cap 1–5) | XP (Cap²) or Debt Clock |
| A one-time ship passage | Nothing (if routine) or a String/Debt Clock (if dramatic) |
| A magical item (single use) | 2–4 XP or barter/String |
| A permanent magical item | 6–10 XP + quest |
| A favor from an NPC | String (if already earned) or Debt Clock (if new) |
| To clear a Debt Clock segment | 2 XP or a service scene |

\subsection{SB Menu for Currency Complications}
\label{subsec:currency-sb-menu}

Spend Story Beats (see \ref{sec:story-beats}) to introduce economic complications:

\begin{itemize}
  \item \textbf{1 SB:} A promised barter item is slightly damaged or cursed; the NPC demands a small extra favor.
  \item \textbf{2 SB:} A Debt Clock ticks unexpectedly (a creditor grows impatient). A String is discovered to be counterfeit.
  \item \textbf{3 SB:} A Trust Asset is threatened (rival faction sabotages the safehouse). A major barter deal falls through; the party loses a potential String.
  \item \textbf{4+ SB:} A Trust schism occurs; the party must mediate or lose half their members. A Debt Clock is called in early, at the worst possible moment.
\end{itemize}

\subsection{Sample Trust Advancement Track}
\label{subsec:sample-trust-track}

Present this to players as a menu for spending XP on their Trust:

\begin{center}
\begin{tabular}{ll}
\toprule
\textbf{Investment (XP)} & \textbf{Unlock} \\
\midrule
4 & Increase Trust Debt Clock capacity by 2 \\
6 & Add a Minor Asset (e.g., a contact network in one port) \\
8 & Upgrade an Asset from Minor to Standard \\
10 & Add a Cap 3 Follower (a trusted lieutenant) \\
12 & Gain a permanent Standard String (a faction owes the Trust a favor) \\
15 & Trust becomes a Major Terrestrial Patron (can grant its own Rites) \\
\bottomrule
\end{tabular}
\end{center}

\section{Final Advice}
\label{sec:gm-currency-final}

The currency-free economy is designed to reduce bookkeeping and increase drama. Embrace its flexibility. When in doubt, ask yourself: *“What would make the story more interesting right now?”*

\begin{itemize}
  \item If a player wants something trivial, give it freely.
  \item If they want something significant, ask for XP, a String, or a Debt Clock.
  \item If they want something monumental, make it the focus of an entire session (or arc).
\end{itemize}

Remember: The players’ XP and Strings are not just resources – they are maps of what your players care about. A player who spends XP on a spy network wants intrigue. A player who accepts a Debt Clock for a ship wants nautical adventure. A player who founds a Trust wants to build a legacy. Give them what they paid for.

\begin{tcolorbox}[colback=green!5!white,colframe=green!70!black,title=The Gray Wanderer’s Ledger]
\emph{“I have seen a kingdom bought with a whispered secret and a war started with a broken scale. Coin is for counting what is already dead. XP, Strings, and debts – those are for counting what is still alive. Spend them well.”}
\end{tcolorbox}

